\section{The Jacobian datum}
\label{pa-chap:Jacobian}

Throughout this chapter \(k\) is a field and \(C\) is the challenge curve: a
scheme over \(k\), proper, smooth of relative dimension one, and geometrically
irreducible. The degree-zero relative Picard functor
\(\Pic^0_{C/k} : \Over(\Spec k)^{\op} \to \mathrm{CommGrp}\) was constructed in
\ref{pa-def:pic0Functor}. This chapter introduces the \emph{Jacobian datum} --- a
representing object for \(\Pic^0_{C/k}\) together with the certificates its
construction provides --- and develops what follows formally from it: the
group structure and universal property of the representing object, its
uniqueness up to canonical isomorphism, separatedness (every group scheme over
a field is separated), the homogeneity supplied by translations, and, given a
proper irreducible \emph{Abel source} mapping onto it, properness and geometric
irreducibility.

\section{The representability datum}
\label{pa-sec:jacobian_datum}

\begin{definition}[Jacobian datum]
  \label{pa-def:jacobian_data}
  \lean{AlgebraicGeometry.JacobianData}
  \leanok
  \dcref{custom}

  \uses{pa-def:pic0Functor}
  A \emph{Jacobian datum} for \(C\) consists of:
  \begin{enumerate}
    \item a scheme \(J\) over \(k\);
    \item a representation of the degree-zero relative Picard functor by
      \(J\): a family of bijections
      \(\Hom_{\Over(\Spec k)}(T, J) \cong \mathrm{pic}^0(C, T)\), natural in
      the test object \(T\);
    \item the certificate that the structure morphism of \(J\) is locally of
      finite type;
    \item the certificate that the structure morphism of \(J\) is
      quasi-compact.
  \end{enumerate}
  The two certificates come from the construction of the representing object
  and are carried as part of the datum rather than re-derived.
\end{definition}

\begin{definition}[The group structure of the representing object]
  \label{pa-def:jacobian_data_grp}
  \lean{AlgebraicGeometry.JacobianData.grpObj}
  \leanok
  \dcref{custom}

  \uses{pa-def:jacobian_data, pa-def:pic0Functor}
  For a Jacobian datum \(d\) with representing object \(J\), the scheme \(J\)
  is a group object of \(\Over(\Spec k)\), with multiplication, unit and
  inversion induced by the commutative-group structure of
  \(\mathrm{pic}^0(C, -)\) through the representation.
\end{definition}
\begin{proof}
  \leanok
  Transport the group structure across the natural bijections of
  \ref{pa-def:jacobian_data}: each \(\Hom(T, J)\) becomes a group, naturally in
  \(T\) --- precomposition with any \(T' \to T\) is a group homomorphism
  because the restriction maps of \(\Pic^0_{C/k}\) are. By the Yoneda lemma the
  natural transformation
  \(\Hom(-, J) \times \Hom(-, J) \to \Hom(-, J)\) given by the pointwise
  products is induced by a unique morphism \(\mu : J \otimes J \to J\), and
  similarly for the unit \(1 \to J\) and the inversion \(J \to J\); the group
  axioms hold for these morphisms because they hold pointwise on
  \(T\)-points.
\end{proof}

\begin{lemma}[The universal property, consumer form]
  \label{pa-lem:jacobian_data_homEquiv}
  \lean{AlgebraicGeometry.JacobianData.homEquiv,
    AlgebraicGeometry.JacobianData.homEquiv_comp}
  \leanok
  \dcref{custom}

  \uses{pa-def:jacobian_data, pa-def:pic0Functor}
  Let \(d\) be a Jacobian datum with representing object \(J\). For every test
  object \(T\) there is a bijection
  \[
    \theta_T : \Hom_{\Over(\Spec k)}(T, J) \;\cong\; \mathrm{pic}^0(C, T),
  \]
  natural against restriction: for \(f : T' \to T\) and \(g : T \to J\),
  \[
    \theta_{T'}(g \circ f) \;=\; f^*\bigl(\theta_T(g)\bigr),
  \]
  where \(f^*\) is the restriction map of \(\Pic^0_{C/k}\)
  (\ref{pa-def:pic0Functor}).
\end{lemma}
\begin{proof}
  \leanok
  This is the representation carried by the datum (\ref{pa-def:jacobian_data}),
  with its naturality read element-wise.
\end{proof}

\begin{definition}[Uniqueness of the representing object up to canonical
    isomorphism]
  \label{pa-def:jacobian_data_unique}
  \lean{AlgebraicGeometry.JacobianData.uniqueUpToIso,
    AlgebraicGeometry.JacobianData.homEquiv_uniqueUpToIso_hom}
  \leanok
  \dcref{custom}

  \uses{pa-def:jacobian_data, pa-lem:jacobian_data_homEquiv}
  Let \(d\), \(d'\) be two Jacobian data for the same curve, with representing
  objects \(J\), \(J'\) and universal bijections \(\theta\), \(\theta'\). There
  is a canonical isomorphism \(\sigma : J \cong J'\) in \(\Over(\Spec k)\),
  characterised by the intertwining property
  \[
    \theta'_T\bigl(\sigma \circ f\bigr) = \theta_T(f)
    \qquad \text{for every } f : T \to J .
  \]
\end{definition}
\begin{proof}
  \leanok
  Let \(u = \theta_J(\mathrm{id}_J) \in \mathrm{pic}^0(C, J)\) be the universal
  class of \(d\), and set \(\sigma = \theta'^{-1}_J(u) : J \to J'\).
  For \(f : T \to J\), naturality of \(\theta'\)
  (\ref{pa-lem:jacobian_data_homEquiv}) gives
  \[
    \theta'_T(\sigma \circ f) = f^*\bigl(\theta'_J(\sigma)\bigr)
    = f^*(u) = f^*\bigl(\theta_J(\mathrm{id}_J)\bigr)
    = \theta_T(\mathrm{id}_J \circ f) = \theta_T(f).
  \]
  Symmetrically there is \(\sigma' = \theta^{-1}_{J'}(u') : J' \to J\), where
  \(u' = \theta'_{J'}(\mathrm{id}_{J'})\), with
  \(\theta_T(\sigma' \circ f') = \theta'_T(f')\) for every \(f' : T \to J'\).
  The two composites are identities: e.g.\
  \(\theta_J(\sigma' \circ \sigma) = \theta'_J(\sigma) = u =
  \theta_J(\mathrm{id}_J)\), and \(\theta_J\) is injective, so
  \(\sigma' \circ \sigma = \mathrm{id}_J\); likewise the other composite.
  Uniqueness of an isomorphism with the intertwining property follows from
  injectivity of \(\theta'\) at \(T = J\).
\end{proof}

\section{Functoriality of the datum in the curve}
\label{pa-sec:jacobian_datum_functoriality}

Throughout this section \(X\), \(Y\) and \(Z\) are challenge curves over \(k\)
(proper, smooth of relative dimension one, geometrically irreducible), with
Jacobian data \(d_X\), \(d_Y\), \(d_Z\) and representing objects \(J_X\),
\(J_Y\), \(J_Z\).

\begin{theorem}[The Yoneda embedding for group objects]
  \label{pa-thm:yonedaGrp}
  \lean{CategoryTheory.yonedaGrp, CategoryTheory.yonedaGrpFullyFaithful,
    CategoryTheory.yonedaGrpObjIsoOfRepresentableBy}
  \mathlibok
  \leanok
  \dcref{custom}

  Let \(\mathcal C\) be a cartesian monoidal category. Sending a group object
  \(G\) to the presheaf of groups \(\Hom(-, G)\) is a fully faithful embedding
  of the category of group objects of \(\mathcal C\) into presheaves of groups
  on \(\mathcal C\). If a presheaf of groups \(F\) is representable, with
  representing object \(J\), then the transported group structure makes \(J\)
  a group object and \(\Hom(-, J) \cong F\) as presheaves of groups.
\end{theorem}

\begin{definition}[Pullback of degree-zero classes as a morphism of group
    schemes]
  \label{pa-def:jacobian_data_pullbackHom}
  \lean{AlgebraicGeometry.JacobianData.pullbackHom,
    AlgebraicGeometry.JacobianData.yonedaGrp_map_pullbackHom}
  \leanok
  \dcref{custom}

  \uses{pa-def:jacobian_data, pa-def:jacobian_data_grp, pa-def:pic0PullbackNat,
    pa-thm:yonedaGrp}
  Let \(g \colon X \to Y\) be a morphism of curves. The composite natural
  transformation
  \[
    \Hom(-, J_Y) \;\cong\; \Pic^0_{Y/k}
    \xrightarrow{\;g^{*}\;} \Pic^0_{X/k} \;\cong\; \Hom(-, J_X)
  \]
  of presheaves of groups --- the outer isomorphisms being the universal
  properties of the two data, as isomorphisms of presheaves of groups
  (Theorem~\ref{pa-thm:yonedaGrp}), and the middle map the degree-zero curve
  transport (\ref{pa-def:pic0PullbackNat}) --- is induced, by full faithfulness
  of the Yoneda embedding for group objects, by a unique morphism of group
  objects
  \[
    g^{*} \colon J_Y \longrightarrow J_X
  \]
  in \(\Over(\Spec k)\): the pullback of degree-zero line-bundle classes,
  realized on the representing objects.
\end{definition}

\begin{theorem}[Functor laws of the pullback morphism]
  \label{pa-thm:jacobian_data_pullbackHom_laws}
  \lean{AlgebraicGeometry.JacobianData.pullbackHom_id,
    AlgebraicGeometry.JacobianData.pullbackHom_comp}
  \leanok
  \dcref{custom}

  \uses{pa-def:jacobian_data_pullbackHom}
  \(\id_X^{*} = \id_{J_X}\) (both slots taken by the same datum \(d_X\)), and
  for \(g \colon X \to Y\) and \(g' \colon Y \to Z\),
  \[
    (g' \circ g)^{*} \;=\; g^{*} \circ g'^{*} \colon J_Z \longrightarrow J_X .
  \]
\end{theorem}
\begin{proof}
  \leanok
  \uses{pa-def:pic0PullbackNat, pa-thm:yonedaGrp}
  The Yoneda embedding for group objects is faithful
  (Theorem~\ref{pa-thm:yonedaGrp}), so it suffices to compare the conjugated
  transformations of presheaves of groups. For the identity, the middle map is
  the identity transformation (\ref{pa-def:pic0PullbackNat}) and the universal
  isomorphism cancels against its inverse. For the composite, the middle map
  of \((g' \circ g)^{*}\) is the composite of the two middle maps
  (\ref{pa-def:pic0PullbackNat}), and the inner universal isomorphism of \(d_Y\)
  cancels against its inverse between the two conjugates.
\end{proof}

\begin{lemma}[The universal-element characterization]
  \label{pa-lem:jacobian_data_comp_pullbackHom}
  \lean{AlgebraicGeometry.JacobianData.comp_pullbackHom,
    AlgebraicGeometry.JacobianData.homEquiv_comp_pullbackHom}
  \leanok
  \dcref{custom}

  \uses{pa-def:jacobian_data_pullbackHom, pa-lem:jacobian_data_homEquiv,
    pa-def:pic0Pullback}
  Let \(g \colon X \to Y\), let \(T\) be a test object and \(f \colon T \to
  J_Y\). Under the universal bijections \(\theta^X\), \(\theta^Y\) of the two
  data (\ref{pa-lem:jacobian_data_homEquiv}),
  \[
    \theta^X_T\bigl(g^{*} \circ f\bigr)
      \;=\; g^{*}\bigl(\theta^Y_T(f)\bigr),
  \]
  the right-hand transport being the degree-zero pullback of
  \ref{pa-def:pic0Pullback}: on functor points, \(g^{*} \colon J_Y \to J_X\)
  \emph{is} pullback of degree-zero classes.
\end{lemma}
\begin{proof}
  \leanok
  Evaluate the defining natural transformation of
  Definition~\ref{pa-def:jacobian_data_pullbackHom} at \(T\) on the element
  \(f\): the first isomorphism sends \(f\) to \(\theta^Y_T(f)\), the middle
  map transports it to \(g^{*}\bigl(\theta^Y_T(f)\bigr)\), and the last
  isomorphism is \((\theta^X_T)^{-1}\); composing with \(\theta^X_T\) gives
  the display.
\end{proof}

\section{Group schemes are separated}
\label{pa-sec:group_separated}

The diagonal argument below is categorical and uses no finiteness hypothesis.
In particular, it proves a stronger form than the standard statement for group
schemes locally of finite type over a field.

\begin{theorem}[The diagonal of a group object is a pullback of the unit]
  \label{pa-thm:grp_diagonal_pullback}
  \lean{CategoryTheory.GrpObj.isPullback_diagonal}
  \leanok
  \dcref{custom}

  \uses{pa-def:grp_diff}
  Let \(G\) be a group object in a cartesian monoidal category, with unit
  \(e : 1 \to G\) and difference morphism
  \(\delta = \mathrm{pr}_1 \cdot \mathrm{pr}_2^{-1} : G \otimes G \to G\)
  (\ref{pa-def:grp_diff}). Then the square
  \[
    \begin{array}{ccc}
      G & \xrightarrow{\ \Delta = \langle \mathrm{id}, \mathrm{id}\rangle\ }
        & G \otimes G \\
      \downarrow & & \downarrow \delta \\
      1 & \xrightarrow{\ e\ } & G
    \end{array}
  \]
  is a pullback: the diagonal of \(G\) is the base change of the unit along
  the difference map.
\end{theorem}
\begin{proof}
  \leanok
  \emph{Commutativity.} In the hom-group \(\Hom(G, G)\):
  \(\delta \circ \Delta = (\mathrm{pr}_1 \circ \Delta)\cdot(\mathrm{pr}_2
  \circ \Delta)^{-1} = \mathrm{id} \cdot \mathrm{id}^{-1} = 1\), the constant
  unit, which is the composite through the terminal object.

  \emph{Universality.} Let \(s_1 : T \to 1\) and \(s_2 : T \to G \otimes G\)
  satisfy \(\delta \circ s_2 = e \circ s_1\). The right-hand side is the unit
  of \(\Hom(T, G)\), and the left-hand side is
  \((\mathrm{pr}_1 \circ s_2)\cdot(\mathrm{pr}_2 \circ s_2)^{-1}\); hence
  \(\mathrm{pr}_1 \circ s_2 = \mathrm{pr}_2 \circ s_2 =: h\). The morphism
  \(h : T \to G\) satisfies \(\Delta \circ h = \langle h, h \rangle = s_2\)
  (equal components) and \(t_G \circ h = s_1\) (morphisms to the terminal
  object agree). For uniqueness, any \(m\) with \(\Delta \circ m = s_2\) has
  \(m = \mathrm{pr}_1 \circ \Delta \circ m = \mathrm{pr}_1 \circ s_2 = h\).
\end{proof}

\begin{theorem}[A group scheme with closed unit section is separated]
  \label{pa-thm:separated_of_closed_unit}
  \lean{AlgebraicGeometry.isSeparated_of_isClosedImmersion_one}
  \leanok
  \dcref{custom}

  \uses{pa-thm:grp_diagonal_pullback}
  Let \(S\) be a scheme and \(G\) a group object of \(\Over(S)\). If the unit
  section \(e : S \to G\) is a closed immersion of schemes, then the structure
  morphism \(G \to S\) is separated. No finiteness hypothesis is required.
\end{theorem}
\begin{proof}
  \leanok
  Apply the forgetful functor \(\Over(S) \to \mathrm{Sch}\) to the pullback
  square of \ref{pa-thm:grp_diagonal_pullback}; it preserves the square and its
  pullback property (fibre products in \(\Over(S)\) are computed on underlying
  schemes), and the monoidal product \(G \otimes G\) becomes the fibre product
  \(G \times_S G\) with \(\Delta\) the diagonal morphism of \(G\) over \(S\).
  Thus the diagonal \(\Delta_{G/S}\) is a base change of the closed immersion
  \(e\), and closed immersions are stable under base change; a morphism whose
  diagonal is a closed immersion is separated.
\end{proof}

\begin{theorem}[Group schemes over a field are separated]
  \label{pa-thm:group_scheme_separated}
  \lean{AlgebraicGeometry.isSeparated_of_grpObj}
  \leanok
  \dcref{custom}

  \uses{pa-thm:separated_of_closed_unit}
  Every group scheme \(G\) over a field \(K\) is separated over \(K\).
\end{theorem}
\begin{proof}
  \leanok
  The unit section \(e : \Spec K \to G\) of a group scheme over a field is a
  closed immersion: its image is a \(K\)-rational point, and a section of a
  \(K\)-scheme at a rational point is a closed immersion. Apply
  \ref{pa-thm:separated_of_closed_unit}.
\end{proof}

\begin{corollary}[The representing object of a Jacobian datum is separated]
  \label{pa-cor:jacobianData_separated}
  \lean{AlgebraicGeometry.JacobianData.isSeparated}
  \leanok
  \dcref{custom}

  \uses{pa-def:jacobian_data_grp, pa-thm:group_scheme_separated}
  For every Jacobian datum \(d\) with representing object \(J\), the structure
  morphism of \(J\) is separated.
\end{corollary}
\begin{proof}
  \leanok
  \(J\) is a group scheme over \(k\) (\ref{pa-def:jacobian_data_grp}); apply
  \ref{pa-thm:group_scheme_separated}.
\end{proof}

\section{Translations}
\label{pa-sec:translations}

Group schemes are homogeneous: right translation by a point is an
automorphism, so a point-local property holding at one rational point spreads
to every rational point. This section provides the translation isomorphisms
and the transport statements, first for a group object \(G\) in a cartesian
monoidal category \(\mathcal C\) (with the hom-group calculus of
\ref{pa-sec:difference_map}), then on underlying schemes. For a point
\(g : 1 \to G\), the \emph{right translation} \(\rho_g : G \cong G\) is the
isomorphism \(\rho_g = \mu \circ \langle \mathrm{id}_G,\ g \circ t_G\rangle\)
(``multiply on the right by \(g\)''), with inverse the right translation by
\(g^{-1}\).

\begin{lemma}[Composition with a right translation]
  \label{pa-lem:comp_mulRight}
  \lean{CategoryTheory.GrpObj.comp_mulRight_hom,
    CategoryTheory.GrpObj.comp_mulRight_inv}
  \leanok
  \dcref{custom}

  For morphisms \(f : X \to G\) and points \(g : 1 \to G\) of a group object
  \(G\),
  \[
    \rho_g \circ f = f \cdot \hat g
    \qquad\text{and}\qquad
    \rho_g^{-1} \circ f = f \cdot \hat g^{-1},
  \]
  computed in the hom-group \(\Hom(X, G)\), where
  \(\hat g = g \circ t_X\) is the constant morphism at \(g\).
\end{lemma}
\begin{proof}
  \leanok
  Unfolding \(\rho_g\): \(\rho_g \circ f = \mu \circ \langle \mathrm{id},\ g
  \circ t_G\rangle \circ f = \mu \circ \langle f,\ g \circ t_X\rangle
  = f \cdot \hat g\), using \(t_G \circ f = t_X\) (morphisms into the terminal
  object). The second identity follows from the first applied to
  \(\rho_g^{-1} = \rho_{g^{-1}}\), or directly: unfolding the inverse
  translation gives \(\mu \circ \langle f,\ \iota \circ g \circ t_X\rangle =
  f \cdot \hat g^{-1}\).
\end{proof}

\begin{definition}[The translation carrying one point to another]
  \label{pa-def:point_translation}
  \lean{CategoryTheory.GrpObj.pointTranslation}
  \leanok
  \source{kleiman-picard}

  \dcref{custom}
  For points \(x, y : 1 \to G\) of a group object \(G\), the
  \emph{point translation}
  \[
    \tau_{x,y} \;=\; \rho_y \circ \rho_x^{-1} \;:\; G \cong G
  \]
  is the automorphism ``multiply on the right by \(x^{-1} y\)''.
\end{definition}

\begin{lemma}[The point translation carries \(x\) to \(y\)]
  \label{pa-lem:point_translation_apply}
  \lean{CategoryTheory.GrpObj.comp_pointTranslation_hom}
  \leanok
  \dcref{custom}

  \uses{pa-def:point_translation, pa-lem:comp_mulRight}
  For points \(x, y : 1 \to G\),
  \(\tau_{x,y} \circ x = y\).
\end{lemma}
\begin{proof}
  \leanok
  By \ref{pa-lem:comp_mulRight}, computed in the hom-group \(\Hom(1, G)\), where
  the constant \(\hat g\) at a point \(g\) is \(g\) itself (\(t_1 =
  \mathrm{id}\)):
  \(\rho_x^{-1} \circ x = x \cdot x^{-1} = 1\), and then
  \(\rho_y \circ (\rho_x^{-1} \circ x) = 1 \cdot y = y\).
\end{proof}

\begin{definition}[Translations of the underlying scheme]
  \label{pa-def:point_translation_iso}
  \lean{CategoryTheory.GrpObj.pointTranslationIso,
    CategoryTheory.GrpObj.pointTranslationIso_hom,
    CategoryTheory.GrpObj.pointTranslationIso_inv}
  \leanok
  \dcref{custom}

  \uses{pa-def:point_translation}
  Let \(S\) be a scheme, \(G\) a group scheme over \(S\), and \(x, y\) two
  sections \(S \to G\). The \emph{translation isomorphism}
  \(\tau_{x,y} : G \cong G\) of underlying schemes is the image of the point
  translation \ref{pa-def:point_translation} under the forgetful functor from
  schemes over \(S\) to schemes; its inverse is the image of the inverse point
  translation.
\end{definition}

\begin{lemma}[Translations commute with the structure morphism]
  \label{pa-lem:point_translation_iso_over}
  \lean{CategoryTheory.GrpObj.pointTranslationIso_hom_comp}
  \leanok
  \dcref{custom}

  \uses{pa-def:point_translation_iso}
  In the situation of \ref{pa-def:point_translation_iso}, composing
  \(\tau_{x,y}\) with the structure morphism \(G \to S\) gives the structure
  morphism: \(\tau_{x,y}\) is an automorphism of \(G\) over \(S\).
\end{lemma}
\begin{proof}
  \leanok
  The point translation is a morphism in the category of schemes over \(S\),
  and the forgetful functor remembers the commutation with the structure
  morphisms.
\end{proof}

\begin{lemma}[Translations act on points as prescribed]
  \label{pa-lem:point_translation_iso_points}
  \lean{CategoryTheory.GrpObj.pointTranslationIso_hom_apply}
  \leanok
  \dcref{custom}

  \uses{pa-def:point_translation_iso, pa-lem:point_translation_apply}
  In the situation of \ref{pa-def:point_translation_iso}, for every point
  \(s \in S\),
  \[
    \tau_{x,y}\bigl(x(s)\bigr) = y(s)
  \]
  on underlying topological spaces.
\end{lemma}
\begin{proof}
  \leanok
  The equation \(\tau_{x,y} \circ x = y\) of \ref{pa-lem:point_translation_apply}
  holds in \(\Over(S)\), hence on underlying schemes, hence on underlying
  topological spaces; evaluate at \(s\).
\end{proof}

\begin{lemma}[Transport of the smooth locus]
  \label{pa-lem:smooth_locus_transport}
  \lean{AlgebraicGeometry.mem_smoothLocus_iff_of_comp_eq}
  \leanok
  \dcref{custom}

  Let \(f : X \to S\) be a morphism locally of finite presentation and
  \(e : X \to X\) an open immersion with \(f \circ e = f\). Then for every
  \(z \in X\),
  \[
    e(z) \in \operatorname{SmoothLocus}(f)
    \iff
    z \in \operatorname{SmoothLocus}(f).
  \]
\end{lemma}
\begin{proof}
  \leanok
  Smoothness of \(f\) at a point is local on \(X\): composing with the open
  immersion \(e\), the morphism \(f\) is smooth at \(e(z)\) if and only if
  \(f \circ e\) is smooth at \(z\) (an open immersion is smooth, identifies
  local rings, and smoothness is a stalk-local property of the morphism). Now
  rewrite \(f \circ e = f\).
\end{proof}

\begin{lemma}[Transport of stalk reducedness]
  \label{pa-lem:stalk_reduced_transport}
  \lean{AlgebraicGeometry.isReduced_stalk_iff_of_isOpenImmersion}
  \leanok
  \dcref{custom}

  Let \(e : X \to Y\) be an open immersion of schemes and \(z \in X\). Then
  \(\struct{X,z}\) is reduced if and only if \(\struct{Y, e(z)}\) is reduced.
\end{lemma}
\begin{proof}
  \leanok
  An open immersion induces an isomorphism on stalks, and reducedness is
  invariant under ring isomorphism.
\end{proof}

\begin{lemma}[Transport of irreducibility]
  \label{pa-lem:irreducible_transport}
  \lean{AlgebraicGeometry.isIrreducible_preimage_iff_of_isIso}
  \leanok
  \dcref{custom}

  Let \(e : X \to Y\) be an isomorphism of schemes and \(t \subseteq Y\) a
  subset. Then \(e^{-1}(t)\) is irreducible if and only if \(t\) is
  irreducible.
\end{lemma}
\begin{proof}
  \leanok
  An isomorphism of schemes is a homeomorphism of underlying spaces;
  irreducibility is a topological property, preserved and reflected by
  homeomorphisms (the image of an irreducible set under a continuous map is
  irreducible, applied to \(e\) and \(e^{-1}\)).
\end{proof}

\begin{corollary}[Homogeneity of the representing object]
  \label{pa-cor:jacobian_translation}
  \lean{AlgebraicGeometry.JacobianData.pointTranslation,
    AlgebraicGeometry.JacobianData.comp_pointTranslation_hom,
    AlgebraicGeometry.JacobianData.pointTranslationIso,
    AlgebraicGeometry.JacobianData.pointTranslationIso_hom_comp,
    AlgebraicGeometry.JacobianData.pointTranslationIso_hom_apply,
    AlgebraicGeometry.JacobianData.pointTranslationIso_mem_smoothLocus_iff,
    AlgebraicGeometry.JacobianData.isReduced_stalk_pointTranslationIso_iff,
    AlgebraicGeometry.JacobianData.isIrreducible_pointTranslationIso_preimage_iff}
  \leanok
  \dcref{custom}

  \uses{pa-def:jacobian_data_grp, pa-def:point_translation,
    pa-lem:point_translation_apply, pa-def:point_translation_iso,
    pa-lem:point_translation_iso_over, pa-lem:point_translation_iso_points,
    pa-lem:smooth_locus_transport, pa-lem:stalk_reduced_transport,
    pa-lem:irreducible_transport}
  Let \(d\) be a Jacobian datum with representing object \(J\), and let
  \(x, y : \Spec k \to J\) be two \(k\)-points. With the group structure of
  \ref{pa-def:jacobian_data_grp}, the translation \(\tau_{x,y}\) is an
  automorphism of \(J\) over \(k\) with \(\tau_{x,y} \circ x = y\); on the
  underlying scheme it commutes with the structure morphism and carries the
  point of \(x\) to the point of \(y\); and it preserves and reflects
  membership in the smooth locus of the structure morphism (when the latter is
  locally of finite presentation), reducedness of stalks, and irreducibility
  of subsets.
\end{corollary}
\begin{proof}
  \leanok
  Instantiate \ref{pa-def:point_translation}--\ref{pa-lem:irreducible_transport} at
  the group scheme \(J\): the translation exists and carries \(x\) to \(y\)
  (\ref{pa-def:point_translation}, \ref{pa-lem:point_translation_apply}); its
  underlying isomorphism commutes with the structure morphism
  (\ref{pa-lem:point_translation_iso_over}) and acts on points as prescribed
  (\ref{pa-lem:point_translation_iso_points}); the three transport lemmas
  (\ref{pa-lem:smooth_locus_transport}, applicable by
  \ref{pa-lem:point_translation_iso_over}; \ref{pa-lem:stalk_reduced_transport};
  \ref{pa-lem:irreducible_transport}) give the invariances.
\end{proof}

\section{The Abel source: properness and geometric irreducibility}
\label{pa-sec:abel_source}

\begin{definition}[Abel source]
  \label{pa-def:abel_source}
  \lean{AlgebraicGeometry.AbelSourceData}
  \leanok
  \dcref{custom}

  \uses{pa-def:jacobian_data}
  Let \(d\) be a Jacobian datum with representing object \(J\). An
  \emph{Abel source} for \(d\) consists of:
  \begin{enumerate}
    \item a scheme \(D\) over \(k\);
    \item a certificate that \(D\) is proper over \(k\);
    \item a certificate that \(D\) is geometrically irreducible over \(k\);
    \item an \emph{Abel morphism} \(a : D \to J\) over \(k\);
    \item a certificate that \(a\) is surjective on underlying topological
      spaces.
  \end{enumerate}
  This is a conditional interface: its definition asserts neither the
  existence of \(D\) nor the existence of \(a\). The intended instance takes
  \(D\) to be a proper scheme of effective divisors of large degree on \(C\)
  and \(a\) to be the Abel map.
\end{definition}

\begin{theorem}[The representing object is universally closed]
  \label{pa-thm:jacobianData_universallyClosed}
  \lean{AlgebraicGeometry.universallyClosed_of_abelSource}
  \leanok
  \dcref{custom}

  \uses{pa-def:abel_source, pa-def:jacobian_data}
  Let \(d\) be a Jacobian datum admitting an Abel source. Then the structure
  morphism of the representing object \(J\) is universally closed.
\end{theorem}
\begin{proof}
  \leanok
  Let \(a : D \to J\) be the Abel morphism. Since \(a\) is a morphism over
  \(k\), the composite of \(a\) with the structure morphism of \(J\) is the
  structure morphism of \(D\), which is universally closed (\(D\) is proper).
  Surjectivity and closedness are stable under base change, so it suffices to
  see: if \(g \circ f\) is universally closed and \(f\) is surjective, then
  \(g\) is universally closed. Indeed, after any base change, for a closed
  subset \(Z\) of the source of the (base-changed) \(g'\),
  \(g'(Z) = (g' \circ f')\bigl(f'^{-1}(Z)\bigr)\) by surjectivity of the
  base-changed \(f'\), and the right-hand side is closed since
  \(g' \circ f' = (g \circ f)'\) is a closed map.
\end{proof}

\begin{theorem}[The representing object is proper]
  \label{pa-thm:jacobianData_proper}
  \lean{AlgebraicGeometry.isProper_of_abelSource}
  \leanok
  \dcref{custom}

  \uses{pa-def:abel_source, pa-cor:jacobianData_separated,
    pa-thm:jacobianData_universallyClosed, pa-def:jacobian_data}
  Let \(d\) be a Jacobian datum admitting an Abel source. Then the structure
  morphism of the representing object \(J\) is proper.
\end{theorem}
\begin{proof}
  \leanok
  Properness is the conjunction of separatedness, universal closedness and
  finite type. Separatedness is \ref{pa-cor:jacobianData_separated}; universal
  closedness is \ref{pa-thm:jacobianData_universallyClosed}; the structure morphism
  is locally of finite type by the datum's certificate
  (\ref{pa-def:jacobian_data}), and quasi-compact because universally closed
  morphisms are quasi-compact; locally of finite type and quasi-compact is
  finite type.
\end{proof}

\begin{theorem}[The representing object is geometrically irreducible]
  \label{pa-thm:jacobianData_geomIrreducible}
  \lean{AlgebraicGeometry.geometricallyIrreducible_of_abelSource}
  \leanok
  \dcref{custom}

  \uses{pa-def:abel_source}
  Let \(d\) be a Jacobian datum admitting an Abel source. Then the structure
  morphism of the representing object \(J\) is geometrically irreducible: for
  every field extension \(K/k\), the base change \(J_K\) is irreducible.
\end{theorem}
\begin{proof}
  \leanok
  Fix a field extension \(K/k\) and a fibre product \(Z = J \times_k K\) with
  projections to \(J\) and \(\Spec K\). The base change
  \(D_K = D \times_k K\) is irreducible by the Abel source's certificate. The
  comparison morphism \(D_K \to Z\) induced by the Abel morphism \(a\) is,
  by pasting of pullback squares, a base change of \(a\) along
  \(Z \to J\) (the square with legs \(D_K \to Z \to \Spec K\) over
  \(D \to J \to \Spec k\) decomposes into two pullbacks); surjectivity is
  stable under base change, so \(D_K \to Z\) is surjective. A continuous
  surjective image of an irreducible space is irreducible, so \(Z\) is
  irreducible.
\end{proof}

\section{Base change of the representability datum}
\label{pa-sec:jacobian_datum_basechange}

Fix a field extension \(k \to L\), write \(\sigma : \Spec L \to \Spec k\) and
\(C_L = C \times_{\Spec k} \Spec L\). We transport a Jacobian datum for \(C\) to one for
\(C_L\) and compare the two group structures on the base-changed representing object; this is the
datum-level mechanism behind \ref{pa-thm:baseChangeIso}.

\begin{definition}[Base change of a Jacobian datum]
  \label{pa-def:jacobianData_baseChange}
  \lean{AlgebraicGeometry.JacobianData.baseChange}
  \leanok
  \dcref{custom}

  \uses{pa-def:jacobian_data, pa-def:baseChange, pa-lem:jacobian_data_homEquiv}
  A Jacobian datum \(d\) for \(C\), with representing object \(J\), base-changes to a Jacobian datum
  \(d_L\) for \(C_L\) with representing object \(J_L = J \times_{\Spec k} \Spec L\). Its universal
  bijections are obtained from the pullback--base-change adjunction
  \(\Hom_{\Over(\Spec L)}(T, J_L) \cong \Hom_{\Over(\Spec k)}(\sigma_! T, J)\) followed by \(d\)'s
  representation and the base-field comparison of the degree-zero Picard functor
  \(\mathrm{pic}^0(C_L, T) \cong \mathrm{pic}^0(C, \sigma_! T)\); the locally-of-finite-type and
  quasi-compactness certificates transport by stability of these properties under base change.
\end{definition}

\begin{lemma}[Base change respects the group structure]
  \label{pa-lem:grpObjObj_baseChange_eq}
  \lean{AlgebraicGeometry.grpObjObj_baseChange_eq}
  \leanok
  \dcref{custom}

  \uses{pa-def:jacobianData_baseChange, pa-def:jacobian_data_grp}
  Let \(d\) be a Jacobian datum for \(C\). The group-object structure that base change induces on
  \(J_L\) --- as the image of the group object \(J\) under the finite-product-preserving base-change
  functor --- coincides with the group-object structure carried by the base-changed datum \(d_L\)
  (\ref{pa-def:jacobianData_baseChange}) through its representation.
\end{lemma}
\begin{proof}
  \leanok
  A group-object structure is determined by its multiplication morphism, so it suffices to compare
  those. The multiplication induced by a finite-product-preserving functor on the image of a group
  object is, through the pullback--base-change adjunction, exactly the multiplication that the
  adjunction representation \(\Hom(T, J_L) \cong \Hom(\sigma_! T, J)\) transports from \(J\). The
  representation carried by \(d_L\) factors that same adjunction through the base-field comparison of
  \(\mathrm{pic}^0\), which is an isomorphism of presheaves of \emph{groups}; being a homomorphism at
  each test object it preserves the pointwise products, so the two multiplications agree.
\end{proof}

\begin{definition}[Datum-level base-change isomorphism]
  \label{pa-def:baseChangeIsoOfData}
  \lean{AlgebraicGeometry.baseChangeIsoOfData}
  \leanok
  \dcref{custom}

  \uses{pa-lem:grpObjObj_baseChange_eq, pa-def:jacobianData_baseChange, pa-def:jacobian_data_unique}
  Let \(d\) be a Jacobian datum for \(C\) and \(d_L\) a Jacobian datum for \(C_L\). There is a
  canonical isomorphism of group schemes over \(L\)
  \[
    J_L \;\cong\; d_L.J ,
  \]
  where the left-hand side carries the group structure induced by base change. It is the canonical
  isomorphism of representing objects (\ref{pa-def:jacobian_data_unique}) for the two Jacobian data of
  \(C_L\) --- the base-changed datum of \ref{pa-def:jacobianData_baseChange} and the given \(d_L\) ---
  which is a homomorphism of group objects by \ref{pa-lem:grpObjObj_baseChange_eq} together with that
  intertwining property. Specialised to the produced Jacobian data it is the datum-level form of
  \ref{pa-thm:baseChangeIso}.
\end{definition}
